```latex
\documentclass[11pt,a4paper]{article}

\usepackage[utf8]{inputenc}
\usepackage[T1]{fontenc}
\usepackage{lmodern}
\usepackage{amsmath,amssymb,amsfonts}
\usepackage{graphicx}
\usepackage{booktabs}
\usepackage{geometry}
\usepackage{hyperref}
\usepackage{physics}
\usepackage{bm}
\usepackage{authblk}

\geometry{margin=2.5cm}

\title{
\textbf{Paper 21 --- The SLACS Fixed Point}\\
\large Strong-Lensing Validation of the BuP Effective Potential
}

\author{Farid Hamdad}
\affil{Bottom-Up Quantum Gravity Program}
\date{2026}

\begin{document}

\maketitle

\begin{abstract}
We test a prediction of the Bottom-Up Quantum Gravity framework on the SLACS strong-lensing sample. Paper 9 introduced an effective gravitational potential generated by an emergent matter source through a graph Poisson equation,
\[
L_{\rm ent}\Phi_{\rm BuP}=S_{\rm flux}.
\]
Here we test whether this potential predicts strong-lensing residuals beyond a baryonic reference model. The observable is
\[
C_{\rm obs}
=
\log\left(
\frac{\theta_E^{\rm obs}}
{\theta_E^{\rm baryon}}
\right).
\]
We find a robust transition around
\[
\log M_\star\simeq 11.58-11.60,
\]
where \(C_{\rm obs}\simeq0\), where the BuP potential locally outperforms the global stellar-mass proxy \(\log M_\star\), and where the dimensional sector satisfies \(\alpha_{\rm eff}\simeq1\) at an intermediate diffusion scale. Using measured Sersic indices strengthens the signal relative to a forced \(n=4\) control. These results provide a non-trivial observational consistency test of the BuP effective potential predicted in Paper 9.
\end{abstract}

\section{Introduction}

The Bottom-Up Quantum Gravity program proposes that geometry and gravity emerge from the structure of quantum entanglement. Previous papers developed the theoretical chain
\[
W_{ij}
\rightarrow
L_{\rm ent}
\rightarrow
d_s,d_w
\rightarrow
\alpha_{\rm eff}
\rightarrow
\text{effective gravity}.
\]

Paper 9 introduced an emergent matter source \(S_{\rm flux}\) and coupled it to an effective graph Poisson equation,
\[
L_{\rm ent}\Phi_{\rm BuP}=S_{\rm flux}.
\]
The present paper tests this prediction on the SLACS strong-lensing sample.

The central question is whether \(\Phi_{\rm BuP}\) leaves an observable imprint in lensing residuals.

\section{Prediction from Paper 9}

Paper 9 predicts the chain
\[
S_{\rm flux}
\rightarrow
\Phi_{\rm BuP}
\rightarrow
\text{gravitational response}.
\]

The source is
\[
S_{\rm flux}
=
T_{00}
-\frac{1}{2}T_{aa}
+\frac{1}{2}T_{\rm grad}
+
|T_{0a}|.
\]

The effective potential is obtained through
\[
L_{\rm ent}\Phi_{\rm BuP}=S_{\rm flux}.
\]

In the lensing application, the BuP potential is tested through the residual correction
\[
C_{\rm obs}
=
\log\left(
\frac{\theta_E^{\rm obs}}
{\theta_E^{\rm baryon}}
\right).
\]

\section{Dataset and observables}

We use a SLACS strong-lensing working catalogue containing lensing radii, stellar masses, effective radii and photometric profile information.

The baryonic reference prediction is denoted
\[
\theta_E^{\rm baryon}.
\]

The observed Einstein radius is
\[
\theta_E^{\rm obs}.
\]

The residual is
\[
C_{\rm obs}
=
\log\left(
\frac{\theta_E^{\rm obs}}
{\theta_E^{\rm baryon}}
\right).
\]

A positive value indicates an excess over the baryonic reference; a negative value indicates an underprediction or over-normalization depending on the adopted baseline.

\section{BuP graph construction}

For each galaxy, we construct a graph from a projected photometric profile. The baseline parametric surface density is represented by a Sersic profile
\[
\Sigma(R)
\propto
\exp\left[
-b_n
\left(
\frac{R}{R_e}
\right)^{1/n}
\right].
\]

The graph coupling is taken to be
\[
W_{ij}
\propto
\sqrt{\Sigma_i\Sigma_j}
\exp\left(
-\frac{d_{ij}^2}{2\xi^2}
\right).
\]

The entanglement Laplacian is
\[
L_{\rm ent}=D-W.
\]

The potential is obtained from the pseudo-inverse structure of the graph Laplacian and from the Paper 9 source construction.

\section{Dimensional sector}

The effective gravitational exponent is
\[
\alpha_{\rm eff}
=
\frac{2d_s}{d_w}+d_w-4.
\]

The fixed-point condition is
\[
\alpha_{\rm eff}=1.
\]

Equivalently,
\[
2d_s=d_w(5-d_w),
\]
or
\[
d_s=\frac{d_w(5-d_w)}{2}.
\]

For standard Brownian diffusion,
\[
d_w=2,
\]
and the fixed point gives
\[
d_s=3.
\]

Thus \(\alpha_{\rm eff}=1\) corresponds to the Newtonian or baryonic fixed point of the effective gravitational sector.

\section{Global residual test}

We first compare the BuP potential to the baryonic residual
\[
C_{\rm obs}.
\]

The best dynamical BuP feature reaches a leave-one-out improvement of approximately \(10.49\%\) on the full working sample.

This establishes that \(\Phi_{\rm BuP}\) carries predictive information about the lensing residual.

\begin{table}[h]
\centering
\begin{tabular}{lcc}
\toprule
Model & LOO improvement & Interpretation \\
\midrule
Baseline baryonic & 0\% & Reference \\
LOW/HIGH regime & \(\sim 6-8\%\) & Discrete BuP regime signal \\
Shape-only BuP force & \(\sim 9\%\) & Morphological graph signal \\
\(\Phi_{\rm BuP}^{\sqrt{M_\star}}\) & \(10.49\%\) & Dynamical BuP potential \\
\(\sqrt{M_\star}\) & \(12.25\%\) & Stellar-mass proxy \\
\(\log M_\star\) & \(12.92\%\) & Best global one-parameter proxy \\
\bottomrule
\end{tabular}
\caption{Summary of the residual-prediction hierarchy.}
\end{table}

\section{Dynamic shuffle control}

To test whether the signal is caused by the correct association between stellar amplitude and graph structure, we perform a dynamic shuffle control.

The stellar amplitude is permuted between galaxies before constructing the graph.

This destroys most of the BuP signal, reducing the best shuffled potential to a small residual contribution.

This demonstrates that the result depends on the correct amplitude-structure association, not merely on the marginal distribution of stellar masses.

\section{Measured Sersic indices}

A controlled comparison is performed on the same 61 galaxies with measured Sersic indices.

Using measured \(n_i\), the best dynamical BuP feature reaches
\[
10.40\%
\]
LOO improvement.

For the exact same galaxies forced to \(n=4\), the best dynamical BuP feature drops to
\[
9.12\%.
\]

Thus, the measured photometric morphology strengthens the BuP dynamical signal.

\begin{table}[h]
\centering
\begin{tabular}{lcc}
\toprule
Catalogue & \(N\) & Best dynamical BuP LOO improvement \\
\midrule
Measured \(n_i\) & 61 & \(10.40\%\) \\
Forced \(n=4\) & 61 & \(9.12\%\) \\
Measured + fallback & 70 & \(10.49\%\) \\
\bottomrule
\end{tabular}
\caption{Effect of using measured Sersic indices.}
\end{table}

\section{Transition window}

A sliding-window analysis in stellar mass shows that the BuP potential locally outperforms \(\log M_\star\) around
\[
11.545<\log M_\star<11.645.
\]

With measured Sersic indices, in this window,
\[
\Phi_{\rm BuP}
\]
outperforms \(\log M_\star\) for approximately
\[
81.8\%
\]
of galaxies.

This suggests that the BuP potential captures a localized dynamical correction in the transition regime.

\section{Observation fixed point}

The observational fixed point is defined by
\[
C_{\rm obs}=0.
\]

Fitting \(C_{\rm obs}\) as a function of \(\log M_\star\) gives a zero crossing around
\[
\log M_\star\simeq 11.58-11.60.
\]

This coincides with the transition window in which \(\Phi_{\rm BuP}\) outperforms the global stellar-mass proxy.

\section{Dimensional fixed point}

A scan over diffusion windows shows that the dimensional fixed point
\[
\alpha_{\rm eff}=1
\]
is recovered at an intermediate diffusion scale.

For the measured-Sersic sample, the optimal window is
\[
t_{\min}=1,\qquad t_{\max}=21.
\]

This gives
\[
\langle \alpha_{\rm eff}\rangle=1.014,
\]
with
\[
d_w\simeq4.213,
\qquad
d_s\simeq1.688.
\]

Inside the transition mass window,
\[
11.545<\log M_\star<11.645,
\]
the same diffusion window gives
\[
\langle \alpha_{\rm eff}\rangle=1.014.
\]

Thus the dimensional fixed point and the observational fixed point coincide.

\section{Triple convergence}

The main result of Paper 21 is the triple convergence
\[
C_{\rm obs}=0,
\qquad
\Phi_{\rm BuP}>\log M_\star,
\qquad
\alpha_{\rm eff}\simeq1,
\]
around
\[
\log M_\star\simeq11.6.
\]

This is the SLACS fixed point.

It is not introduced as a free threshold. It emerges independently from the residual lensing data, from the BuP potential performance window, and from the graph-dimensional sector.

\section{Discussion}

The result should not be interpreted as a complete validation of the entire BuP program. It is a validation of a specific prediction of Paper 9: the effective BuP potential carries observable gravitational information.

The current graph remains parametric. A stronger test would use non-parametric HST light maps to build a genuinely two-dimensional graph.

Nevertheless, the fact that measured Sersic indices strengthen the result indicates that the BuP potential is sensitive to real photometric morphology, not only to stellar mass.

\section{Conclusion}

We have tested the Paper 9 BuP effective potential on the SLACS strong-lensing sample.

The main findings are:

\begin{enumerate}
\item The BuP potential predicts lensing residuals with a non-zero LOO improvement.
\item Dynamic shuffle controls suppress the signal.
\item Measured Sersic indices strengthen the BuP prediction.
\item A transition window appears around \(\log M_\star\simeq11.6\).
\item The observed residual satisfies \(C_{\rm obs}=0\) near the same mass.
\item The dimensional sector reaches \(\alpha_{\rm eff}\simeq1\) at an intermediate diffusion scale in the same mass window.
\end{enumerate}

Therefore, Paper 21 provides an observational validation of the Paper 9 prediction that the BuP effective potential can encode gravitational lensing corrections.

\bibliographystyle{unsrt}
\bibliography{references}

\end{document}
```
